\documentclass{article}
\usepackage{amsmath, amssymb, amsthm}
\usepackage{hyperref}
\usepackage{geometry}
\geometry{margin=1in}

\title{Erd\H{o}s Problem \#945}
\date{Last updated: 2025-08-31}
\author{Source: erdosproblems.com}

\begin{document}
\maketitle

\noindent\textbf{Status:} OPEN \quad \textbf{No prize}

\noindent\textbf{Tags:} number theory, divisors

\medskip

\noindent\textbf{Problem Statement:}

Let $F(x)$ be the maximal $k$ such that there exist $n+1,\ldots,n+k\leq x$ with $\tau(n+1),\ldots,\tau(n+k)$ all distinct (where $\tau(m)$ counts the divisors of $m$). Estimate $F(x)$. In particular, is it true that\[F(x) \leq (\log x)^{O(1)}?\]In other words, is there a constant $C&#62;0$ such that, for all large $x$, every interval $[x,x+(\log x)^C]$ contains two integers with the same number of divisors?




A problem of Erd\H{o}s and Mirsky \cite{ErMi52}, who proved that\[\frac{(\log x)^{1/2}}{\log\log x}\ll F(x) \ll \exp\left(O\left(\frac{(\log x)^{1/2}}{\log\log x}\right)\right).\]Erd\H{o}s \cite{Er85e} claimed that the lower bound could be improved via their method 'with some more work' to $(\log x)^{1-o(1)}$. Beker has improved the upper bound to\[F(x) \ll \exp\left(O\left((\log x)^{1/3+o(1)}\right)\right).\]Cambie has observed that Cram\'{e r's conjecture} implies that $F(x) \ll (\log x)^2$, and furthermore if every interval in $[x,2x]$ of length $\gg \log x$ contains a squarefree number (see [208] ) then every interval of length $\gg (\log x)^2$ contains two numbers with the same number of divisors, whence\[F(x) \ll (\log x)^2.\]See [1004] for the analogous problem with the Euler totient function. This problem is discussed in problem B18 of Guy's collection \cite{Gu04}.




References

[Er85e] Erd\H{o}s, P., Some problems and results in number theory . Number theory and combinatorics. Japan 1984 (Tokyo, Okayama and Kyoto, 1984) (1985), 65-87.

[ErMi52] Erd\H{o}s, P. and Mirsky, L., The distribution of values of the divisor function {$d(n)$} . Proc. London Math. Soc. (3) (1952), 257--271.

[Gu04] Guy, Richard K., Unsolved problems in number theory . (2004), xviii+437.

\noindent\textbf{Related OEIS sequences:} possible, A048892


\bigskip
\noindent\small{Source: \url{https://www.erdosproblems.com/945}}
\end{document}
